\documentclass[11pt]{article}
\usepackage[a4paper,margin=1in]{geometry}
\usepackage{amsmath,amssymb}
\usepackage{booktabs}
\usepackage{array}

\title{MinScorer Report}
\author{DifficultyAgri}
\date{\today}

\begin{document}
\maketitle

\section{Purpose}
The MinScorer estimates how difficult each training image is for object detection. It measures difficulty at two levels:
\begin{itemize}
	\item \textbf{Object level}: how difficult each ground-truth object is to detect.
	\item \textbf{Image level}: how difficult the full image is after aggregating object difficulty and false positives.
\end{itemize}
The resulting scores are used later for dataset analysis, object mining, and copy-paste augmentation.

\section{Inputs}
The scorer consumes:
\begin{itemize}
	\item Ground-truth labels in YOLO format.
	\item Low-confidence predictions from the model.
	\item Optimal-confidence predictions from the model.
	\item A scoring configuration with parameters \(\alpha\), \(\beta\), IoU threshold, object weight, false-positive weight, and weight mode.
\end{itemize}

\section{Object Difficulty}
For each ground-truth object, the scorer searches for predictions with the same class and IoU above the threshold. The object difficulty is defined as
\[
S_{obj} = \min_{p} \left[ \alpha(1-\mathrm{conf}(p)) + \beta(1-\mathrm{IoU}(g,p)) \right].
\]
If no valid match is found, the score defaults to the maximum penalty
\[
S_{obj} = \alpha + \beta.
\]

This means:
\begin{itemize}
	\item higher confidence lowers the score,
	\item higher IoU lowers the score,
	\item unmatched objects are treated as the hardest cases.
\end{itemize}

\section{Image Difficulty}
For each image, the scorer computes the average object difficulty:
\[
\overline{S}_{obj} = \frac{1}{N} \sum_{i=1}^{N} S_{obj}^{(i)}.
\]
It also computes the false-positive rate and the missed-detection rate using the optimal-confidence predictions.

The final image score is
\[
S_{img} = w_1 \cdot \overline{S}_{obj} + w_2 \cdot FP_{rate}.
\]

\section{How Each Parameter Is Determined}
\subsection*{\(\alpha\)}
User-defined. It controls how strongly low confidence increases object difficulty.

\subsection*{\(\beta\)}
User-defined. It controls how strongly poor localization increases object difficulty.

\subsection*{IoU threshold}
User-defined. A prediction is considered a valid match only if its IoU with the ground-truth object is at least this value.

\subsection*{Object weight \(w_1\)}
User-defined. It scales the contribution of average object difficulty in the image score.

\subsection*{False-positive weight \(w_2\)}
Determined by the selected \texttt{weight\_mode}:
\begin{itemize}
	\item \textbf{fixed}: use the configured false-positive weight directly.
	\item \textbf{mean\_match}: set \(w_2\) so that the average false-positive contribution matches the average object-score scale.
	\item \textbf{balance\_correlation}: search for a value of \(w_2\) that balances score correlation with false positives and missed detections.
\end{itemize}

\section{Current Configuration Values}
The current experiment configuration uses the following values:

\begin{center}
\begin{tabular}{>{\raggedright\arraybackslash}p{0.45\linewidth} p{0.35\linewidth}}
\toprule
Parameter & Value \\
\midrule
Scorer type & \texttt{min\_scorer} \\
\(\alpha\) & 0.5 \\
\(\beta\) & 0.5 \\
IoU threshold & 0.5 \\
Object weight \(w_1\) & 0.9 \\
False-positive weight & 0.1 \\
Weight mode & \texttt{balance\_correlation} \\
\bottomrule
\end{tabular}
\end{center}

\section{Interpretation}
\begin{itemize}
	\item A higher object score means the object is harder to detect.
	\item A higher image score means the image is harder overall.
	\item The false-positive rate raises the image score even when object matches are strong.
\end{itemize}

\section{Summary}
The MinScorer turns detector predictions into a structured difficulty score. It combines object-level matching quality and image-level false-positive behavior into a single score that can be used for ranking samples and guiding augmentation.

\end{document}
